\section{Introduction}
\label{sec:introduction}

A retinal vessel segmentation model that averages 0.78 Dice across a test set
tells a clinician very little about any individual image. Some of those images
were segmented near-perfectly; others contain errors that, if acted upon, could
affect screening decisions for diabetic retinopathy or glaucoma. The aggregate
metric hides this variance. What a deployment actually needs is not a better
average score, but a policy: for each image (or each pixel), should the system's
output be trusted, flagged for review, or handed off to a human?

This question sits at the intersection of uncertainty estimation and selective
prediction, two areas that have developed largely in parallel within medical
image analysis. Uncertainty estimation methods, from Bayesian approximations
\citep{gal2016dropout} to ensemble techniques \citep{lakshminarayanan2017simple},
produce maps that correlate (to varying degrees) with model error. Selective
prediction, studied more thoroughly in classification
\citep{geifman2017selective,el2010foundations}, formalizes the idea that a model
should abstain when its confidence is low, trading coverage for accuracy.
In segmentation, the connection between these two ideas remains
underdeveloped. Most work stops at generating uncertainty maps and computing
calibration metrics, without asking how those maps should drive downstream
decisions.

We address this gap directly. Rather than proposing a new segmentation
architecture or a new uncertainty estimator, we study the \emph{decision layer}
that sits on top of existing methods. Given an uncertainty map and a prediction,
what policy should a system follow to decide which outputs to accept and which
to defer? The answer matters more than it might seem: the right combination of
uncertainty method and deferral policy eliminates roughly 80\% of segmentation
errors while deferring only 25\% of pixels. Uncertainty estimation alone,
without a decision mechanism to act on it, leaves this reduction on the table.

This framing changes what counts as a good uncertainty method. A well calibrated
model is not necessarily useful if its uncertainty fails to separate correct
pixels from incorrect ones. A poorly calibrated model whose uncertainty
reliably flags errors may work fine for deferral. We test this distinction
and find that it matters: temperature scaling barely moves calibration metrics
while sometimes degrading the uncertainty-error correlation that deferral
actually depends on.

Our study is built on retinal vessel segmentation using the DRIVE dataset, with
cross-dataset evaluation on STARE and CHASE\_DB1. We use a U-Net with a
ResNet-34 encoder as the base segmentation model and compare two uncertainty
estimation approaches: MC Dropout (30 stochastic forward passes) and Test-Time
Augmentation (6 geometric transforms). On top of these, we evaluate three
deferral strategies of increasing sophistication:

\begin{enumerate}[leftmargin=*]
  \item \textbf{Global thresholding}: defer all pixels whose uncertainty exceeds
  a fixed threshold, chosen to optimize a deferral F1 score on a validation set.
  \item \textbf{Image-adaptive thresholding}: compute a per-image threshold
  based on percentile statistics of that image's uncertainty distribution, so
  the deferral rate adapts to each image's difficulty.
  \item \textbf{Confidence-aware deferral}: weight uncertainty by the inverse of
  prediction confidence, so that uncertain pixels near the decision boundary
  ($p \approx 0.5$) are deferred preferentially over uncertain pixels where the
  model is nonetheless confident.
\end{enumerate}

We evaluate these combinations using risk-coverage analysis, deferral curves,
and a direct measure of error reduction as a function of deferral rate. We also
test temperature scaling as a post-hoc calibration step and evaluate its effect
on both calibration metrics (ECE) and decision quality (deferral error
reduction).

Our contributions:

\begin{enumerate}[leftmargin=*]
  \item We show that TTA uncertainty discriminates correct from incorrect pixels
  far better than MC Dropout in the unified comparison bundle
  (Unc-AUROC 0.881 vs.\ 0.722) while running about 3.1$\times$ faster,
  making it the preferred uncertainty source for any
  deferral-based pipeline on this task.

  \item We propose and compare three deferral policies of increasing
  sophistication. The best overall configuration, TTA with adaptive deferral,
  reduces error by approximately 80\% at 25\% deferral. A confidence-aware
  variant, $s = u \cdot (1 - c)$, achieves the best \emph{efficiency}: 55\%
  error reduction at just 12\% deferral while using a smaller review budget
  than the global TTA policy.

  \item We demonstrate that temperature scaling, the standard post-hoc
  calibration fix, does not improve the tracked deferral operating points in
  this repository. In the main comparison bundle it worsens the reported ECE
  values for both MC Dropout and TTA, illustrating that calibration metrics
  and deferral utility should be evaluated separately.

  \item We provide cross-dataset evidence (DRIVE $\to$ STARE, CHASE\_DB1)
  that uncertainty quality transfers more robustly than segmentation accuracy
  under domain shift: the MC Dropout pipeline retains Unc-AUROC above 0.83 on
  both external datasets despite a roughly 6.5-point Dice drop from DRIVE.

  \item We introduce a risk-coverage evaluation framework with clinically
  motivated operating points (screening, diagnostic, aggressive) that lets
  practitioners select deferral behavior based on their specific review
  budget and error tolerance.
\end{enumerate}

Sections~\ref{sec:related_work}--\ref{sec:experimental_setup} cover related
work, our methodology, and experimental setup.
Sections~\ref{sec:results}--\ref{sec:ablations} present claim-driven results
and ablations. Sections~\ref{sec:discussion}--\ref{sec:conclusion} discuss
implications, limitations, and conclusions.
